
\documentclass{article}
\usepackage{colortbl}
\usepackage{makecell}
\usepackage{multirow}
\usepackage{supertabular}

\begin{document}

\newcounter{utterance}

\centering \large Interaction Transcript for game `cladder', experiment `full\_v1.5\_default', episode 5155 with refineonly{-}Qwen3.5{-}4B.
\vspace{24pt}

{ \footnotesize  \setcounter{utterance}{1}
\setlength{\tabcolsep}{0pt}
\begin{supertabular}{c@{$\;$}|p{.15\linewidth}@{}p{.15\linewidth}p{.15\linewidth}p{.15\linewidth}p{.15\linewidth}p{.15\linewidth}}
    \# & \multicolumn{2}{c}{Player} && \multicolumn{2}{c}{Game Master} \\
    \hline

    \theutterance \stepcounter{utterance}  
    & & & \multicolumn{4}{p{0.6\linewidth}}{
        \cellcolor[rgb]{0.9,0.9,0.9}{
            \makecell[{{p{\linewidth}}}]{
                \texttt{\tiny{[P1$\langle$GM]}}
                \texttt{You are an expert in causal inference. The following question is not a typical commonsense query, but rather a meticulously designed question created by a professor specializing in causal inference, intended to assess the students' mastery of the course content.} \\
\texttt{Imagine a self{-}contained, hypothetical world with only the following conditions, and without any unmentioned factors or causal relationships: Talent has a direct effect on elite institution admission status. Brown eyes has a direct effect on elite institution admission status. Method 1: We look at how talent correlates with brown eyes case by case according to elite institution admission status. Method 2: We look directly at how talent correlates with brown eyes in general. To understand how talent affects brown eyes, is it more correct to use the Method 1 than Method 2?} \\
\texttt{Start your answer with 'Yes' or 'No', followed by additional reasoning or evidence to support your explanation.} \\
            }
        }
    }
     \\ \\

    \theutterance \stepcounter{utterance}  
    & \multicolumn{4}{p{0.6\linewidth}}{
        \cellcolor[rgb]{0.561,0.737,0.561}{
            \makecell[{{p{\linewidth}}}]{
                \texttt{\tiny{[P1$\rangle$GM]}}
                \texttt{No} \\
\\ 
\texttt{**Reasoning:**} \\
\\ 
\texttt{In this hypothetical world, we are asked to understand the causal effect of **Talent** on **Brown Eyes**. However, the problem statement explicitly defines the causal structure as:} \\
\texttt{1.  Talent \$\textbackslash{}rightarrow\$ Elite Institution Admission} \\
\texttt{2.  Brown Eyes \$\textbackslash{}rightarrow\$ Elite Institution Admission} \\
\\ 
\texttt{There is **no direct causal link** stated between Talent and Brown Eyes. In causal inference, asking "how Talent affects Brown Eyes" when there is no direct arrow between them usually implies looking for an indirect effect (mediated through the admission status) or determining if they are independent.} \\
\\ 
\texttt{Let's analyze the two methods:} \\
\\ 
\texttt{*   **Method 2 (General Correlation):** This calculates the raw correlation between Talent and Brown Eyes across the entire population. Since both variables independently cause a third variable (Admission), they are likely to be positively correlated due to **selection bias** (or collider bias) at the "Admission" node. If high Talent leads to admission and Brown Eyes leads to admission, people with both traits are overrepresented in the admitted group, creating a spurious positive correlation even if Talent has *no* causal effect on Brown Eyes. This method captures the association, not the causal effect.} \\
\\ 
\texttt{*   **Method 1 (Case{-}by{-}case according to admission status):** This phrasing implies stratifying the analysis by the admission status (e.g., comparing admitted vs. non{-}admitted groups, or looking at the relationship within strata).} \\
\texttt{    *   If we stratify by **Admission Status** (the collider), we induce **selection bias**. By conditioning on the common effect (Admission), we create a spurious correlation between the parents (Talent and Brown Eyes) that did not exist causally. This would make it *worse* to estimate any causal relationship.} \\
\texttt{    *   However, if the question implies using the admission status as a **control variable** (adjusting for it) to isolate the relationship, we must be careful. In a DAG where \$Talent \textbackslash{}rightarrow Admission \textbackslash{}leftarrow BrownEyes\$, adjusting for the collider (\$Admission\$) opens a backdoor path and creates a false association.} \\
\\ 
\texttt{**The Core Logical Flaw in the Premise:**} \\
\texttt{The question asks, "To understand how talent affects brown eyes..."} \\
\texttt{Since there is **no causal path** from Talent to Brown Eyes in the defined world, the true causal effect is **zero**.} \\
\\ 
\texttt{*   **Method 2** will likely yield a non{-}zero correlation (due to the collider bias at Admission), leading one to falsely conclude that Talent affects Brown Eyes.} \\
\texttt{*   **Method 1** (if interpreted as stratifying on the collider) will also yield a non{-}zero correlation (often stronger or in a different direction depending on the distribution), reinforcing the false conclusion.} \\
\\ 
\texttt{Therefore, neither method correctly identifies the causal effect of zero. However, the question asks which is *more* correct} \\
            }
        }
    }
    & & \\ \\

    \theutterance \stepcounter{utterance}  
    & & & \multicolumn{2}{p{0.3\linewidth}}{
        \cellcolor[rgb]{0.9,0.9,0.9}{
            \makecell[{{p{\linewidth}}}]{
                \texttt{\tiny{[GM$|$GM]}}
                \texttt{No} \\
            }
        }
    }
    & & \\ \\

    \theutterance \stepcounter{utterance}  
    & & & \multicolumn{2}{p{0.3\linewidth}}{
        \cellcolor[rgb]{0.9,0.9,0.9}{
            \makecell[{{p{\linewidth}}}]{
                \texttt{\tiny{[GM$|$GM]}}
                \texttt{no} \\
            }
        }
    }
    & & \\ \\

    \theutterance \stepcounter{utterance}  
    & & & \multicolumn{2}{p{0.3\linewidth}}{
        \cellcolor[rgb]{0.9,0.9,0.9}{
            \makecell[{{p{\linewidth}}}]{
                \texttt{\tiny{[GM$|$GM]}}
                \texttt{game\_result = WIN} \\
            }
        }
    }
    & & \\ \\

\end{supertabular}
}

\end{document}
